\begin{figure*}[t]
\centering
\includegraphics[width=\linewidth]{figure/framework.pdf}
%\caption{fig1}
\centering
% \captionsetup{font={small}}
% \vspace{-1.0em}
\caption{Our \ours architecture comprises three integral parts in memory storage. During note construction, the system processes new interaction memories and stores them as notes with multiple attributes. The link generation process first retrieves the most relevant historical memories and then employs an LLM to determine whether connections should be established between them. The concept of a 'box' describes that related memories become interconnected through their similar contextual descriptions, analogous to the Zettelkasten method. However, our approach allows individual memories to exist simultaneously within multiple different boxes. During the memory retrieval stage, we extract query embeddings using a text encoding model and search the memory database for relevant matches. When related memory is retrieved, similar memories that are linked within the same box are also automatically accessed.}
\label{fig:framework}
\vspace{-1.0em}
\end{figure*}
